\section{Исследовательская часть}
\label{sec:research}

Раздел вводит базовые понятия Рефала-5 и разбирает, как объектные выражения хранят
существующие реализации языка. Сперва~--- минимум терминологии, на которую опираются
дальнейшие рассуждения: программа, функция, предложение, объектное выражение,
сопоставление с образцом. Затем~--- сами представления и их компромиссы, а в итоге
обоснование того, почему для проектируемой системы выбрана раскрученная двусторонняя
очередь.

\subsection{Программа, функции и предложения}
\label{sec:refal-program}

Программа на Рефале~--- это набор функций. Каждая функция задаётся одним или
несколькими предложениями (правилами) вида \enq{образец = результат},
заключёнными в фигурные скобки. При вызове функции система просматривает её
предложения сверху вниз и берёт первое, чей образец сопоставим с
аргументом~\mbox{\cite{refal_book,refal05_site}}. По подстановке, полученной при
сопоставлении, строится результирующее выражение, которое заменяет вызов в
текущем поле зрения.

Простейший пример~--- рекурсивная функция, вычисляющая длину объектного
выражения (листинг~\ref{lst:length}):

\begin{lstlisting}[language=Refal,caption={Функция подсчёта длины выражения},label={lst:length}]
Length {
   = 0;
  s.First e.Rest = <Add 1 <Length e.Rest>>;
}
\end{lstlisting}

Функция \prog{Length} имеет два предложения. Первое срабатывает на пустом
выражении и возвращает \prog{0}. Второе сопоставляется с непустым выражением:
оно связывает первый терм с переменной \prog{s.First}, остаток~--- с
переменной \prog{e.Rest} и рекурсивно подсчитывает длину остатка.
Имена с префиксами \prog{s.}, \prog{t.}, \prog{e.} обозначают переменные
образца разных типов; их семантика рассматривается ниже.

\subsection{Объектные выражения, термы и символы}

Все данные в Рефале представлены объектными выражениями.
Объектное выражение~--- это конечная последовательность термов
(возможно пустая). Каждый терм имеет один из двух видов:

\begin{itemize}
  \item атомарный терм~--- собственно символ в терминологии Рефала;
        в Рефале-5 различают три разновидности символов: символ-слово
        (идентификатор: \prog{Found}, \prog{True}, \prog{Nil}),
        макроцифра (целое число произвольной длины) и литера (один знак
        строкового литерала);
  \item скобочный терм~--- объектное выражение, заключённое в круглые скобки.
\end{itemize}

Важно не смешивать понятия \enq{символ} и \enq{литера}. Символ-слово~---
это атомарный объект (метка), а не последовательность литер: при
сопоставлении \prog{Found} ведёт себя как один атомарный терм. Запись в
одинарных кавычках, например \prog{'abc'}, обозначает три литеры подряд,
то есть объектное выражение длины три. Аналогично, число \prog{42}~--- это
один символ (макроцифра), а не последовательность из двух цифр-литер.

Сопоставление с образцом опирается на три класса переменных образца:

\begin{itemize}
  \item \(s\)-переменная (\prog{s.X}) связывается с одним
        символом~--- одним атомарным термом любой из трёх
        разновидностей;
  \item \(t\)-переменная (\prog{t.X}) связывается с одним
        термом, в том числе со скобочным;
  \item \(e\)-переменная (\prog{e.X}) связывается с произвольной,
        возможно пустой, цепочкой термов.
\end{itemize}

Простой пример с переменной \(e\)-типа~--- функция, возвращающая последний
символ непустой последовательности атомарных термов (листинг~\ref{lst:last}):

\begin{lstlisting}[language=Refal,caption={Функция получения последнего символа},label={lst:last}]
Last {
  e.Init s.Last = s.Last;
}
\end{lstlisting}

При вызове \prog{<Last 1 2 3>} переменная \prog{e.Init} связывается с
выражением \prog{1 2}, а \prog{s.Last}~--- с символом \prog{3}. Вариант, в
котором \prog{e.Init} забирала бы всё выражение целиком, не годится:
тогда для \prog{s.Last} не осталось бы ни одного терма, и сопоставление
было бы неудачным. Алгоритм сопоставления учитывает требования к концам
образца и подбирает разбиение, при котором все переменные связываются
согласованно.

\subsection{Неоднозначное сопоставление}
\label{sec:ambiguous-matching}

Если в образце встречается несколько \(e\)-переменных, сопоставление
становится неоднозначным: одно и то же объектное выражение можно
разбить разными способами. Стандартный пример~--- функция, проверяющая
вхождение заданного символа в выражение (листинг~\ref{lst:contains}):

\begin{lstlisting}[language=Refal,caption={Неоднозначное сопоставление при поиске символа},label={lst:contains}]
Contains {
  s.Needle e.Left s.Needle e.Right = True;
  s.Needle e.Any                   = False;
}
\end{lstlisting}

Первое предложение требует, чтобы и в начале, и где-то дальше в выражении
встретился один и тот же символ \prog{s.Needle}. Образец содержит две
\(e\)-переменные \prog{e.Left} и \prog{e.Right}, и при сопоставлении
система перебирает варианты: она начинает с пустого \prog{e.Left}, и
если очередной \prog{s.Needle} не оказывается на текущей позиции,
удлиняет \prog{e.Left} на один терм и пробует снова. Каждое такое
удлинение~--- это комбинация операций \enq{взять один терм с начала
остатка} и \enq{дописать этот терм в конец накапливаемого префикса}, то
есть \emph{отделение начала} и \emph{конкатенация}.

Так даже логически простое правило разворачивается в множество шагов, и на каждом
интерпретатор разбирает и собирает фрагменты выражения. Сколько стоит один такой шаг,
решает представление выражения в памяти~--- к нему и переходим.

\subsection{Объектное выражение как абстрактный тип}
\label{sec:adt-ops}

Чтобы говорить о стоимости разных представлений, удобно зафиксировать набор
операций, через которые выражается интерпретация любого предложения:

\begin{itemize}
  \item проверка на пустоту;
  \item отделение начала: получить первый терм и остаток выражения;
  \item отделение конца: получить последний терм и остаток выражения;
  \item конкатенация двух выражений~--- основная операция при
        построении результата правила;
  \item выражение из одного терма (создание сингулярного выражения).
\end{itemize}

Эти операции вызываются на каждом шаге сопоставления и каждом шаге сборки
результата. Например, при переборе \(e\)-переменной
(\ref{sec:ambiguous-matching}) последовательно работают отделение начала
и конкатенация; при возврате результата правила вида \prog{= <F s.X> <G
e.Y>;} итоговое выражение собирается из нескольких фрагментов с помощью
конкатенации. Поэтому выбор представления~--- это, по существу, выбор
того, какие из перечисленных операций должны быть дёшевы, а какие могут
стать дорогими.

Дальнейший разбор идёт по нарастанию компромисса. Списковые представления дёшевы на
концах и в конкатенации, но платят \emph{копированием} значений переменных. Массивные и
векторно-списковые меняют беду местами: копирование дешевеет, зато дорожает
\emph{конкатенация}. Дерево конкатенаций показано как гипотетический промежуточный
вариант. Раскрученная двусторонняя очередь, которой раздел заканчивается, амортизированно
снимает обе проблемы сразу.

\subsection{Классическое (плоское двусвязное) представление}
\label{sec:flat-list}

Этот способ представления использовался в большинстве реализаций Рефала
вплоть до 1980-х годов и часто называется классическим~\cite{skorobogatov_repr}
или плоским двусвязным~\cite{abramov_romanenko_1988}. Объектное
выражение хранится как двусвязный список звеньев, где каждое звено
соответствует либо атомарному символу, либо одной из двух структурных
скобок. У звена-скобки есть дополнительное поле, указывающее на парную
скобку. Контейнер выражения хранит указатели на первое и последнее звенья.

Стоимость основных операций при таком представлении:

\begin{itemize}
  \item отделение начала или конца --- \(O(1)\): сдвинуть указатель; если
        крайнее звено --- открывающая скобка, можно прыгнуть через поле
        парной скобки на закрывающую и получить скобочный терм целиком;
  \item конкатенация двух выражений --- \(O(1)\): достаточно перенастроить
        четыре указателя на стыке;
  \item копирование выражения --- \(O(N)\), где \(N\) --- число
        звеньев в записи (включая все вложенные скобки): необходимо
        физически создать новый список звеньев.
\end{itemize}

Дешёвая конкатенация делает это представление удобным для
построения результата, но за неё приходится платить копированием: если переменная
встречается в правой части предложения чаще, чем в левой, каждое \enq{лишнее} вхождение
требует полной копии её значения. Стандартная иллюстрация --- функция
подстановки из препринта Абрамова и Романенко~\cite{abramov_romanenko_1988}:
для каждого символа выражения по таблице ищется его замена.

\begin{lstlisting}[language=Refal,caption={Подстановка по таблице (пример проблемы копирования)},label={lst:subst-table}]
Subst {
  (e.Tab)            = ;
  (e.Tab) s.X e.Rest = <Lookup (e.Tab) s.X> <Subst (e.Tab) e.Rest>;
}
\end{lstlisting}

Переменная \prog{e.Tab} в левой части появляется один раз, а в правой ---
дважды (передаётся в \prog{Lookup} и в рекурсивный вызов \prog{Subst}).
При классическом представлении это означает, что на каждом шаге рекурсии
таблица копируется целиком. Если длина исходного выражения и длина
таблицы оба порядка \(n\), общая стоимость становится \(O(n^2)\), и
попытка программировать в функциональном стиле упирается в эту
квадратичную асимптоту. Распространённый \enq{инженерный} обход --- так
называемый \enq{сквозной просмотр выражения} с двумя стеками, переводящий
\prog{Subst} в императивный по духу вариант без многократного
использования \prog{e.Tab}; копирований он избегает ценой более тёмного кода и
потерянной возможности распараллеливания~\cite{abramov_romanenko_1988}.

\subsection{Компромиссное представление с подвешенными скобками}
\label{sec:hanging-list}

Чтобы смягчить проблему копирования, не отказываясь от списковой основы,
было предложено компромиссное представление~\cite{skorobogatov_repr}.
Оно используется в реализации Рефала-6 и ряде учебных проектов, а также
обсуждается в документации Рефала-05~\cite{refal05_site}.

В этом представлении звено списка соответствует не отдельному знаку
записи, а одному терму верхнего уровня. Скобочный терм хранится
особым образом: вместо двух парных звеньев (открывающей и закрывающей
скобок) звено содержит ссылку на отдельный кольцевой список, представляющий
вложенное выражение. У головы кольца хранится счётчик ссылок.

Стоимость операций по сравнению с плоским списком:

\begin{itemize}
  \item отделение начала и конца --- \(O(1)\), как и раньше;
  \item конкатенация --- \(O(1)\);
  \item копирование верхнего уровня --- \(O(N_{\text{верх}})\),
        где \(N_{\text{верх}}\) --- число термов верхнего уровня;
        вложенные выражения при копировании \emph{не дублируются}, у их
        кольцевых списков лишь увеличивается счётчик ссылок.
\end{itemize}

Применительно к функции \prog{Subst} это даёт ощутимый выигрыш в
типичных сценариях: если таблица состоит из пар \prog{(s.Key e.Value)},
то при копировании \prog{e.Tab} физически дублируются только звенья
верхнего уровня, а вложенные пары остаются разделёнными между копиями.

Платой становится более сложное управление памятью: необходимо
поддерживать счётчики ссылок при разрушении и сборке выражений и аккуратно
обрабатывать ситуации, когда переменная связана со \emph{внутренностью}
скобочного терма, который оказался расшарен (тогда модификация одной
\enq{копии} могла бы повредить другие). Кроме того, наличие циклических
связей у голов колец усложняет автоматическую сборку мусора.

\subsection{Массивное представление}
\label{sec:array-repr}

Альтернативный путь, предложенный в работе Абрамова и
Романенко~\cite{abramov_romanenko_1988}, состоит в том, чтобы хранить
объектное выражение в непрерывном участке памяти (массиве
тегированных ячеек). Каждая ячейка содержит либо атомарный символ, либо
открывающую или закрывающую структурную скобку. Скобочный терм при таком
хранении не требует отдельного указателя: он занимает диапазон ячеек от
открывающей до закрывающей скобки.

Главное преимущество массивного представления --- дешёвое копирование
переменной. Значение \(e\)-переменной \prog{e.X} представляется парой
индексов \((l, r)\) внутри массива; копирование переменной --- это просто
копирование пары чисел, то есть \(O(1)\). Это превращает функциональный
стиль программирования в Рефале из \enq{дорогого} в \enq{нормальный}: даже
многократное использование одной и той же переменной в правой части не
приводит к физическому дублированию данных. На функции \prog{Subst} из
\ref{sec:flat-list} массивное представление даёт линейную асимптотику без
каких-либо ухищрений с императивным переписыванием.

Стоимость основных операций:

\begin{itemize}
  \item отделение начала и конца --- \(O(1)\): сдвинуть один из индексов;
  \item копирование переменной --- \(O(1)\);
  \item конкатенация двух выражений длины \(|X|\) и \(|Y|\) ---
        \(O(|X| + |Y|)\): требуется выделить новый массив подходящей
        длины и физически скопировать в него содержимое обоих фрагментов.
\end{itemize}

Так возникает проблема конкатенации, симметричная проблеме
копирования. Если результат правила собирается в цикле, как в типовой
функции отображения

\Needspace{10\baselineskip}
\begin{lstlisting}[language=Refal,caption={Построение результата через рекурсивную конкатенацию},label={lst:map}]
Map {
   = ;
  s.X e.Rest = <F s.X> <Map e.Rest>;
}
\end{lstlisting}

то на \(n\) шагах рекурсии возникают \(n\) конкатенаций. При наивной
массивной реализации каждая из них копирует уже накопленный префикс, и
суммарная стоимость становится \(O(n^2)\).

В препринте~\cite{abramov_romanenko_1988} предлагается смягчение --- так
называемые дырки (запас свободного места) по краям массива. Если
к выражению с дыркой справа дописывается короткий фрагмент длины \(k\),
то вместо создания нового массива достаточно скопировать только эти \(k\)
ячеек в свободную область. На сериях коротких приписываний это даёт
амортизированно линейное общее время. Однако если дырка исчерпывается,
приходится снова выделять новый массив и переносить всё содержимое;
гарантию \(O(1)\) на произвольных последовательностях операций дырки не
дают. Дополнительно в~\cite{abramov_romanenko_1988} обсуждаются меры по
снижению затрат на конкатенацию --- распознавание частных случаев,
функции с несколькими входами и выходами, использование скобок для
структуризации данных.

\subsection{Векторно-списковое представление}
\label{sec:vector-list}

Векторно-списковое представление, предложенное
Скоробогатовым~\cite{skorobogatov_repr}, нацелено на то, чтобы сохранить
дешёвое копирование переменных от массивного представления и одновременно
избежать линейной конкатенации.

Идея состоит в том, чтобы изменить интерфейс между функциями. Аргументы
функций передаются как векторы (массивы тегированных ячеек), но
результат функция возвращает не как один сплошной вектор, а как
список ссылок на фрагменты векторов. Конкатенация фрагментов
откладывается: пока вычисление не дошло до точки, где результат должен
стать пассивным аргументом другой функции (в частности, попасть внутрь
скобочного терма), фрагменты остаются отдельными звеньями.

Применительно к функции \prog{Map} это означает следующее. На каждом шаге
рекурсии вместо физического присоединения \prog{<F s.X>} к результату
\prog{<Map e.Rest>} формируется лёгкое звено \enq{сначала такой-то
фрагмент, потом такой-то}. Когда вычисление функции завершено и её
результат должен стать пассивным выражением, накопленный список
фрагментов собирается в единый вектор за один проход, что даёт линейное
общее время вместо квадратичного.

Платой становится более сложный интерфейс языка сборки и более
ответственное управление временем жизни фрагментов: пока на фрагмент есть
ссылки в незавершённых вычислениях, исходный вектор нельзя освободить или
переиспользовать. Подробное описание представления и связанной системы
команд абстрактной Рефал-машины приведено в~\cite{skorobogatov_repr}.

\subsection{Гипотетическое представление: дерево конкатенаций}
\label{sec:rope-repr}

Среди представлений, не получивших распространения в реализациях Рефала,
интерес представляет дерево конкатенаций. В этом
представлении объектное выражение хранится в виде двоичного дерева, где
листья --- небольшие фрагменты массива, а внутренние узлы соответствуют
операциям конкатенации поддеревьев.

Преимущество rope в том, что конкатенация двух выражений сводится к
созданию одного нового внутреннего узла, что даёт \(O(1)\) с периодической
перебалансировкой за \(O(\log n)\). Копирование переменной также \(O(1)\)
--- это ссылка на корень поддерева. Слабая сторона --- операции на
концах: чтобы получить первый или последний терм, нужно спуститься по
дереву к крайнему листу, что стоит \(O(\log n)\) и зависит от баланса.
Для Рефала это особенно неудобно, потому что \(e\)-переменные требуют как
раз многократного отделения термов с краёв.

В литературе по реализациям Рефала rope обсуждается лишь как теоретическая
возможность~\cite{refal05_site}.

\subsection{Предлагаемое представление: раскрученная двусторонняя очередь}
\label{sec:bootstrapped-deque}

У каждого из разобранных представлений есть своё слабое место:
плоский список --- проблема копирования; компромиссный --- неполное её
устранение и сложное управление памятью; массивное --- проблема
конкатенации; векторно-списковое --- сложный интерфейс и привязка к
определённой модели вычисления; rope --- дорогие операции на концах.
В работе вместо них взята раскрученная двусторонняя
очередь из монографии
Окасаки~\cite{okasaki_pfds}: она амортизированно закрывает обе
ключевые проблемы разом.

Идея структуры опирается на два уровня вложенности.

\begin{itemize}
  \item Мелкий уровень (Shallow) --- буфер фиксированной
        ёмкости \(B\), организованный как циклический массив. Локальные
        операции на концах при наличии места выполняются как сдвиг
        индексов, без выделения памяти.
  \item Глубокий уровень (Deep) --- тройка \(\langle
        front,\, mid,\, back\rangle\), где крайние компоненты
        (\(front\), \(back\)) --- мелкие буферы, а средний (\(mid\)) ---
        двусторонняя
        очередь мелких буферов. Структура определена рекурсивно: средняя
        очередь сама может быть как мелкой, так и глубокой.
\end{itemize}

Структуру держит один инвариант: внутренние буферы средней очереди заполнены не менее
чем наполовину. Без него глубокая очередь могла бы выродиться в длинную
цепочку почти пустых буферов; с ним глубина рекурсивного представления
растёт логарифмически от длины выражения.

Стоимость основных операций:

\begin{itemize}
  \item отделение начала или конца --- \(O(1)\) амортизированно;
  \item конкатенация двух выражений --- \(O(1)\) амортизированно (за счёт
        перетекания крайних буферов в среднюю очередь);
  \item использование переменной несколько раз --- \(O(1)\), благодаря
        структурному разделению: новая версия выражения
        ссылается на те же неизменяемые буферы и узлы, что и старая.
\end{itemize}

Структурное разделение в данном случае --- инженерная техника, а не
свойство языка: в самой модели Рефала-5 нет изменяемых ячеек памяти, и
вычисление концептуально строит новые выражения, а задача реализации в
том, чтобы за этой моделью не стояло физическое копирование.

\subsubsection{Структура мелкого и глубокого уровней}

Рисунок~\ref{fig:shallow_ring} показывает идею мелкого буфера: элементы
размещены циклически, концы задаются индексами \enq{head} и
\enq{tail}; в общем случае head может находиться правее tail, и
разрыв \enq{обходится} по модулю ёмкости \(B\).

\begin{figure}[!htb]
\centering
\begin{tikzpicture}[
    cell/.style={rectangle, draw, minimum width=0.9cm, minimum height=0.55cm, font=\footnotesize, align=center},
    lab/.style={font=\footnotesize},
]
  \node[cell] (c0) { };
  \node[cell, right=0cm of c0] (c1) { };
  \node[cell, right=0cm of c1] (c2) {x};
  \node[cell, right=0cm of c2] (c3) {y};
  \node[cell, right=0cm of c3] (c4) {z};
  \node[cell, right=0cm of c4] (c5) { };
  \node[cell, right=0cm of c5] (c6) { };
  \node[cell, right=0cm of c6] (c7) { };
  \node[cell, right=0cm of c7] (c8) { };
  \node[lab, below=0.25cm of c2] {head};
  \node[lab, below=0.25cm of c4] {tail};
\end{tikzpicture}
\caption{Идея мелкого буфера: элементы размещены циклически, концы задаются индексами}
\label{fig:shallow_ring}
\end{figure}

Глубокий уровень состоит из двух мелких буферов и средней очереди, которая
сама хранит мелкие буферы (рисунок~\ref{fig:deep_deque}). Рекурсивность
означает, что средняя очередь при необходимости тоже становится глубокой,
и мы получаем структуру с логарифмической по основанию \(B/2\) глубиной.

\begin{figure}[!htb]
\centering
\begin{tikzpicture}[
    box/.style={rectangle, draw, rounded corners, minimum width=3.2cm, minimum height=0.85cm, align=center, font=\small},
    arrow/.style={-{Stealth[length=2.5mm]}, thick}
]
    \node[box] (front) {front\\(shallow buffer)};
    \node[box, right=1.2cm of front] (mid) {mid\\(deque of buffers)};
    \node[box, right=1.2cm of mid] (back) {back\\(shallow buffer)};
    \draw[arrow] (front) -- (mid);
    \draw[arrow] (mid) -- (back);
\end{tikzpicture}
\caption{Глубокий уровень раскрученной двусторонней очереди: два мелких буфера и средняя очередь буферов}
\label{fig:deep_deque}
\end{figure}

\subsubsection{Краткое обоснование амортизированных оценок}

Большинство операций затрагивает только крайние буферы \(front\) и
\(back\): пока в них есть свободное место, добавление и удаление
сводятся к сдвигу индексов в циклическом массиве и являются строго
\(O(1)\). Перетекание буфера в среднюю очередь происходит только при
достижении границ заполненности и требует \(O(B)\). Если каждой такой
дорогой операции приписать стоимость, накопленную на предшествующих
локальных операциях, то средняя стоимость шага остаётся около константы;
это и есть классический амортизированный аргумент~\cite{okasaki_pfds}.

Интуитивно редкие глубокие переносы можно сравнить с цепочкой переносов
при прибавлении единицы в позиционной системе счисления: вероятность
цепочки длины \(k\) убывает геометрически, и вклад в \emph{среднюю}
стоимость образует сходящийся ряд. Для основания \(q \ge 2\) ожидаемое
число переносов
\[
  \mathbb{E}[K] = \sum_{k\ge 1} \frac{k(q-1)}{q^{k+1}} = \frac{1}{q-1}.
\]
Это иллюстрация, а не строгое доказательство для конкретной реализации.
Полные техники амортизированного анализа функциональных структур данных
(включая раскрутку и ленивые вычисления) систематически изложены у
Окасаки~\cite{okasaki_pfds}.

\subsubsection{Сводное сравнение представлений}

Сравнение рассмотренных представлений приведено в
таблице~\ref{tab:repr-comparison}. Зафиксирована стоимость трёх ключевых
операций, через которые выражается основная вычислительная нагрузка
исполняющего модуля: отделение терма с конца, конкатенация двух выражений
и использование одного значения переменной несколько раз.

\begin{table}[!htb]
\caption{\centering Сравнение представлений объектных выражений по стоимости основных операций}
\label{tab:repr-comparison}
\centering
\small
\begin{tabularx}{\linewidth}{|>{\raggedright\arraybackslash}p{3.4cm}|>{\centering\arraybackslash}p{1.9cm}|>{\raggedright\arraybackslash}X|>{\raggedright\arraybackslash}p{3.2cm}|}
\hline
\textbf{Представление} & \textbf{Концы} & \textbf{Конкатенация} & \textbf{Копирование переменной} \\
\hline
Классическое (плоское двусвязное) & \(O(1)\) & \(O(1)\) & \(O(N)\) (вся запись) \\
\hline
Компромиссное (подвешенные скобки) & \(O(1)\) & \(O(1)\) & \(O(N_{\text{верх}})\) \\
\hline
Массивное (с дырками) & \(O(1)\) & \(O(|X|+|Y|)\), аморт.~линейная при послед.~приписываниях & \(O(1)\) \\
\hline
Векторно-списковое & \(O(1)\) & откладывается, \(O(1)\) на сборке & \(O(1)\) \\
\hline
Дерево конкатенаций & \(O(\log n)\) & \(O(1)\), перебалансировка \(O(\log n)\) & \(O(1)\) \\
\hline
Раскрученная двусторонняя очередь & \(O(1)\) аморт. & \(O(1)\) аморт. & \(O(1)\) \\
\hline
\end{tabularx}
\end{table}

\subsubsection{Область применимости и ограничения}
\label{subsec:applicability}

Раскрученная двусторонняя очередь оправдана прежде всего тогда, когда
выражения достаточно длинные, а в правилах активно используются
\(e\)-переменные, порождающие многократные разбиения и сборки. На
коротких выражениях и в программах, где почти каждое правило связано с
одним-двумя термами, выигрыш от сложной структуры может уступить простым
спискам или массивам из-за более высокого постоянного множителя.

Целевой реализацией в данной работе является среда JVM. Это естественно
сочетается с неизменяемыми объектами и структурным разделением, но
накладные расходы на работу сборщика мусора нужно учитывать при выборе
ёмкости \(B\) мелкого буфера: слишком малое \(B\) приводит к большому
числу мелких объектов, слишком большое --- ухудшает локальность
модификаций функционального стиля.

\subsection{Выводы по разделу}

Требования к представлению объектных выражений базисного подмножества
Рефала-5 сводятся к следующему: дёшевые операции на концах, дешёвая
конкатенация при построении результата и отсутствие физического
копирования при многократном использовании одной переменной. Существующие
представления закрывают часть этих требований: классическое и
компромиссное списковые --- быстрые на концах и при конкатенации, но
дорогие на копировании; массивное и векторно-списковое --- быстрые на
копировании, но имеют ограничения по конкатенации или по интерфейсу
языка сборки; rope --- быстрая конкатенация и копирование, но дорогие
операции на концах. Раскрученная двусторонняя очередь покрывает все три
требования амортизированно за \(O(1)\) и поэтому выбрана в качестве
представления объектных выражений в проектируемой системе.
